\section{Proactive Secret Resharing}
\label{sec:resharing}

\subsection{Problem statement}

The two-round signing protocol of \S\ref{sec:protocol} is parameterised by a master secret $s \in R_q^{n_{\textrm{vec}}}$, a public matrix $A \in R_q^{m_{\textrm{pk}} \times n_{\textrm{vec}}}$, an error term $e \in R_q^{m_{\textrm{pk}}}$, and the public group key $b = A \cdot s + e$ (rounded into $\tilde b$). At genesis a trusted-dealer Gen distributes Shamir shares of $s$ to the initial validator committee. The secret $s$ never appears explicitly thereafter --- only its shares.

When the validator set rotates at the end of an epoch, the natural na\"{i}ve approach is to run trusted-dealer Gen again with the new committee. This re-derives $s$ inside a single party's memory, which is unacceptable for a leaderless public chain: the new dealer learns the master secret, and the public key $b$ changes (so all signatures from prior epochs become unverifiable without auxiliary indexing).

\textbf{Goal.} Rotate the share distribution from an old committee $\mathcal{O} = \{1, \dots, n_{\mathrm{old}}\}$ with threshold $t_{\mathrm{old}}$ onto a new committee $\mathcal{N}$ (potentially disjoint, with possibly different $|\mathcal{N}|, t_{\mathrm{new}}$), without (i) reconstructing $s$ in any party's memory, and (ii) changing the public key $\tilde b$. Equivalently: any threshold subset of $\mathcal{N}$ jointly produces signatures that verify under the same $\tilde b$ as before.

We adapt the proactive secret-sharing protocol of Herzberg, Jakobsson, Jarecki, Krawczyk, and Yung \cite{hjky97} to the lattice ring $R_q = \mathbb{Z}_q[X]/(X^{256} + 1)$ used by Pulsar.

\subsection{Notation}

Throughout this section, party identifiers are 1-indexed. Each old party $i \in \mathcal{O}$ holds a Shamir share $s_i \in R_q^{n_{\textrm{vec}}}$ of $s$ in standard (non-NTT, non-Montgomery) form, where $s_i$ is the value of an underlying secret polynomial $P : R_q \to R_q^{n_{\textrm{vec}}}$ at $X = i$:
\[
P(X) = s + a_1 X + a_2 X^2 + \cdots + a_{t_{\mathrm{old}} - 1} X^{t_{\mathrm{old}} - 1}, \qquad s_i = P(i),
\]
with $a_1, \dots, a_{t_{\mathrm{old}}-1}$ uniform in $R_q$. Here $P$, $a_d$, and $s_i$ are vectors of polynomials in $R_q^{n_{\textrm{vec}}}$; Shamir is applied independently per coordinate of each component poly. The threshold relation is $\sum_{i \in T} \lambda_i^T \cdot s_i = s$ for any $T \subseteq \mathcal{O}$ with $|T| = t_{\mathrm{old}}$, where the Lagrange basis at $X = 0$ is
\[
\lambda_i^T = \prod_{j \in T,\, j \neq i} \frac{-x_j}{x_i - x_j} \pmod{q}, \qquad x_i = i.
\]

\subsection{Protocol \textsc{Reshare}}

\begin{enumerate}
\item \textbf{Quorum selection.} The protocol picks a deterministic quorum $\mathcal{Q} \subseteq \mathcal{O}$ with $|\mathcal{Q}| = t_{\mathrm{old}}$. The Pulsar canonical takes the smallest-ID quorum so the protocol is reproducible across implementations (load-bearing for the cross-language KAT in \S\ref{sec:resharing-kat}).

\item \textbf{Per-party random polynomial.} Each $i \in \mathcal{Q}$ samples a random degree-$(t_{\mathrm{new}} - 1)$ polynomial $g_i$ in $X$ over $R_q^{n_{\textrm{vec}}}$ with constant term equal to $\lambda_i^{\mathcal{Q}} \cdot s_i$ (the party's own share scaled by its Lagrange coefficient for $\mathcal{Q}$):
\[
g_i(X) = \bigl(\lambda_i^{\mathcal{Q}} \cdot s_i\bigr) + r_{i,1} X + r_{i,2} X^2 + \cdots + r_{i, t_{\mathrm{new}} - 1} X^{t_{\mathrm{new}} - 1},
\]
where $r_{i,d} \in R_q^{n_{\textrm{vec}}}$ is uniform for $1 \le d < t_{\mathrm{new}}$. Each coefficient is sampled independently per Shamir-coordinate slot.

\item \textbf{Share dealing.} For each $j \in \mathcal{N}$, party $i$ computes $g_i(j) \in R_q^{n_{\textrm{vec}}}$ and sends it to party $j$ over an authenticated, encrypted pairwise channel. The transport layer is out of scope for this section --- standard X25519 + ML-KEM hybrid \cite{lp070} suffices.

\item \textbf{New share assembly.} Each new party $j \in \mathcal{N}$ collects the $t_{\mathrm{old}}$ values $\{g_i(j) : i \in \mathcal{Q}\}$ and computes its new share:
\[
s'_j \;=\; \sum_{i \in \mathcal{Q}} g_i(j) \pmod{q}.
\]
\end{enumerate}

\begin{remark}[Comparison with the literature.]
The HJKY97 paper presents two distinct primitives: (a) a same-committee proactive update using a zero-polynomial $z_i(0) = 0$, and (b) a redistribution protocol for changing access structures. The protocol above is a specialization of (b) due to Desmedt and Jajodia \cite{desmedt-jajodia97}; the same-committee version is implemented as a \emph{distinct} routine \textsc{Refresh} (described in \S\ref{sec:resharing-refresh} below) for periodic mobile-adversary defense within a stable validator set.
\end{remark}

\subsection{Protocol \textsc{Refresh} (same-committee)}
\label{sec:resharing-refresh}

\textsc{Refresh} is the kernel for HJKY97's primitive (a). The committee $T = \{\alpha_1, \dots, \alpha_n\}$ and threshold $t$ are held \emph{fixed}. Every party rotates its share to fresh independent randomness while preserving the master secret. This is a proactive-security \emph{hygiene} step within an unchanging validator set --- distinct from \textsc{Reshare} above, which rotates the validator set itself.

\begin{enumerate}
\item Each $i \in T$ samples a random degree-$(t-1)$ polynomial $z_i(X)$ over $R_q^{n_{\textrm{vec}}}$ with $z_i(0) = 0$. Concretely the constant term is forced to zero; the higher-degree coefficients are uniformly random.
\item Each $i$ privately delivers $z_i(\alpha_j)$ to every $j \in T$ over an authenticated channel.
\item Each $j$ updates: $s'_j = s_j + \sum_i z_i(\alpha_j) \pmod q$.
\end{enumerate}

\begin{theorem}[\textsc{Refresh} preserves the master secret]
Define $Z(X) = \sum_i z_i(X)$. $Z$ has degree at most $t-1$ and $Z(0) = 0$. Let $S(X)$ be the original Shamir polynomial of degree $t-1$ with $S(0) = s$ and $S(\alpha_j) = s_j$ for all $j \in T$. The new shares are evaluations of $S'(X) = S(X) + Z(X)$ at the original points; $S'(0) = s$, so the new shares are valid $(t, n)$-Shamir shares of the same master secret $s$. Furthermore the joint distribution of $\{s'_j\}_{j \in T}$ is independent of $\{s_j\}_{j \in T}$ subject to $S'(0) = s$, because $Z$'s degree-$1, \dots, t-1$ coefficients are uniform and independent of $S$.
\end{theorem}

\textsc{Refresh} and \textsc{Reshare} compose: a chain may schedule \textsc{Refresh} on a periodic background timer (e.g.\ every block within a stable epoch) and \textsc{Reshare} only at validator-set rotation boundaries. Both produce activation certificates (\S\ref{sec:resharing-activation}) under the same unchanged group public key.

\subsection{Correctness}

\begin{theorem}[Public-key invariance]
\label{thm:resharing-pk-invariant}
Let $s'_j$ be the share computed by new party $j \in \mathcal{N}$ after \textsc{Reshare}. Then for every threshold subset $T' \subseteq \mathcal{N}$ with $|T'| = t_{\mathrm{new}}$,
\[
\sum_{j \in T'} \lambda_j^{T'} \cdot s'_j \;=\; s \pmod{q},
\]
where $s$ is the master secret of the input shares. Equivalently, the public key $\tilde b$ derived from $s$ at genesis verifies signatures produced by any $t_{\mathrm{new}}$-subset of $\mathcal{N}$.
\end{theorem}

\begin{proof}
Define $G(X) = \sum_{i \in \mathcal{Q}} g_i(X)$. By construction $G$ has degree at most $t_{\mathrm{new}} - 1$ and constant term
\[
G(0) = \sum_{i \in \mathcal{Q}} g_i(0) = \sum_{i \in \mathcal{Q}} \lambda_i^{\mathcal{Q}} \cdot s_i = s,
\]
where the last equality is the Lagrange-interpolation identity for the old quorum. The new shares are $s'_j = G(j)$ for $j \in \mathcal{N}$, so $\{s'_j\}_{j \in \mathcal{N}}$ are valid Shamir shares of $G(0) = s$ with threshold $t_{\mathrm{new}}$. Lagrange interpolation of any $t_{\mathrm{new}}$-subset $T' \subseteq \mathcal{N}$ at $X = 0$ recovers $G(0) = s$. Therefore the underlying master secret is preserved exactly. The genesis-time public key $b = A \cdot s + e$ depends only on $s$ (and the unchanged $A, e$), so $\tilde b = \mathrm{Round}(A \cdot s + e, \xi)$ is identical before and after \textsc{Reshare}.\qed
\end{proof}

\begin{remark}[Degree of $G$]
$G$ has degree exactly $t_{\mathrm{new}} - 1$ except with negligible probability: the leading-coefficient slot of $G$ is $\sum_{i \in \mathcal{Q}} r_{i, t_{\mathrm{new}}-1}$, a sum of $t_{\mathrm{old}}$ independent uniform $R_q^{n_{\textrm{vec}}}$ samples. The probability that this sum equals zero on every coordinate is at most $q^{-n_{\textrm{vec}} \cdot \varphi} = 2^{-48 \cdot 7 \cdot 256} \approx 2^{-86\,016}$, vanishing.
\end{remark}

\subsection{Security}

\subsubsection{Threat model}

We adopt the static-corruption-per-epoch adversary of \cite{hjky97}: in each epoch, a probabilistic polynomial-time adversary $\mathcal{A}$ may corrupt up to $t_{\mathrm{old}} - 1$ parties of $\mathcal{O}$ (resp.\ $t_{\mathrm{new}} - 1$ of $\mathcal{N}$) and learn their shares and per-session randomness. Crucially, corruptions in epoch $e$ and epoch $e+1$ are \emph{independent}: a party corrupted across both epochs is counted against the threshold of each, but the adversary obtains no extra advantage from the temporal aggregation.

We further assume each honest outgoing party's RNG is unbiased and unobservable. A corrupted party's randomness is under adversarial control; a corrupted party may also drop messages, send malformed messages, or otherwise deviate from the protocol --- but the kernel \textsc{Reshare} as specified above does not detect such deviations. Verifiable Secret Resharing (VSR) is layered on top in production; see \S\ref{sec:resharing-vsr}.

\subsubsection{Confidentiality}

\begin{theorem}[Static confidentiality]
\label{thm:resharing-confidentiality}
Let $\mathcal{B}_{\mathrm{old}} \subseteq \mathcal{O}$ with $|\mathcal{B}_{\mathrm{old}}| < t_{\mathrm{old}}$ and $\mathcal{B}_{\mathrm{new}} \subseteq \mathcal{N}$ with $|\mathcal{B}_{\mathrm{new}}| < t_{\mathrm{new}}$ denote the corruption sets in the old and new epochs. Then the joint view of $\mathcal{A}$ across the resharing,
\[
\mathrm{View}_\mathcal{A} = \bigl(\{s_i\}_{i \in \mathcal{B}_{\mathrm{old}}},\, \{s'_j\}_{j \in \mathcal{B}_{\mathrm{new}}},\, \{g_i(j)\}_{i \in \mathcal{Q},\, j \in \mathcal{B}_{\mathrm{new}}}\bigr),
\]
is statistically independent of the master secret $s$.
\end{theorem}

\begin{proof}[Sketch]
The old shares $\{s_i\}_{i \in \mathcal{B}_{\mathrm{old}}}$ alone leak no information about $s$ by the standard $(t_{\mathrm{old}}-1)$-Shamir security, applied per Shamir-coordinate slot \cite{shamir79}. The new shares $\{s'_j\}_{j \in \mathcal{B}_{\mathrm{new}}}$ similarly leak nothing on their own: $|\mathcal{B}_{\mathrm{new}}| < t_{\mathrm{new}}$ so they correspond to $< t_{\mathrm{new}}$ evaluations of $G$. The cross-leakage from the dealing step contributes only $\{g_i(j) : i \in \mathcal{Q},\, j \in \mathcal{B}_{\mathrm{new}}\}$. Each $g_i$ is constructed from $i$'s own share $s_i$ (already known to $\mathcal{A}$ if $i \in \mathcal{B}_{\mathrm{old}}$, otherwise an honest random) plus $t_{\mathrm{new}} - 1$ uniform high-degree coefficients. Conditioned on $s_i$, the distribution of $g_i(j)$ for $|\{j\}| < t_{\mathrm{new}}$ is uniformly random in $R_q^{n_{\textrm{vec}}}$, independent of $s$. Therefore the entire $\mathrm{View}_\mathcal{A}$ admits a simulator that draws the leaked $g_i(j)$ values uniformly without knowing $s$ --- so $\mathcal{A}$'s advantage in distinguishing $s$ from random is information-theoretically zero.\qed
\end{proof}

\begin{remark}[Why we do \emph{not} require $|\mathcal{B}_{\mathrm{old}}| + |\mathcal{B}_{\mathrm{new}}| < t_{\mathrm{old}}$]
A weaker version of \Cref{thm:resharing-confidentiality} would require $|\mathcal{B}_{\mathrm{old}}| + |\mathcal{B}_{\mathrm{new}}| < t_{\mathrm{old}}$ (corruptions must not aggregate across epochs). The stronger statement above holds because every honest old party's contribution to a leaked $g_i(j)$ is randomized by independent $r_{i,d}$ values, and an honest party's $s_i$ is itself uniformly random (per Shamir security). The randomization in degrees $1, \dots, t_{\mathrm{new}} - 1$ is the crucial ingredient: it means the new-epoch corruption set $\mathcal{B}_{\mathrm{new}}$ is "fresh" with respect to $s$, even if the same physical machines are corrupted across both epochs.
\end{remark}

\subsubsection{Public-key indistinguishability}

A separate, weaker security property: an external observer who sees only the public output $\tilde b$ before and after \textsc{Reshare} learns nothing about the rotation. This is trivial: $\tilde b$ is unchanged by \Cref{thm:resharing-pk-invariant}, so there is no observable.

\subsection{Verifiable Secret Resharing (VSR)}
\label{sec:resharing-vsr}

The protocols of \S\ref{sec:resharing}--\S\ref{sec:resharing-refresh} are the arithmetic \emph{kernel}: adequate for honest-but-curious parties and sufficient for cryptographic-correctness proofs. Production deployment in a permissionless setting layers the following components on top, all collectively known as Verifiable Secret Resharing (VSR) \cite{wong02}. Each is implemented in a sibling file in the \texttt{pulsar/reshare/} package.

\paragraph{Authenticated point-to-point channels} (\texttt{pairwise.go}). Each pair $(i, j)$ runs an authenticated key exchange: X25519 in the classical baseline, X25519 $\|$ ML-KEM-768 hybrid in the post-quantum mode. The wire-identity Ed25519 key signs the X25519 ephemeral, defeating active man-in-the-middle. The transcript hash (below) is bound into the signed ephemeral so a stale exchange cannot be replayed in a different invocation.

\paragraph{Domain-separated transcript binding} (\texttt{transcript.go}). Every resharing message and complaint includes the 32-byte BLAKE3 hash of
\[
  (\mathit{chain\_id},\, \mathit{group\_id},\, \mathit{old\_epoch\_id},\, \mathit{new\_epoch\_id},\, H(\mathcal{O}),\, H(\mathcal{N}),\, t_{\mathrm{old}},\, t_{\mathrm{new}},\, H(\tilde b),\, \mathit{variant})
\]
under personalization \texttt{"pulsar.reshare.transcript.v1"}. The \emph{variant} field $\in \{\textsc{refresh}, \textsc{reshare}\}$ separates the two primitives' transcripts so an activation cert produced for one cannot be replayed for the other. The validator-set hashes $H(\mathcal{O}), H(\mathcal{N})$ are computed over sorted public-key byte strings, independent of map iteration order.

\paragraph{Polynomial commitments} (\texttt{commit.go}). Each old party $i$ broadcasts the Pedersen-style commitment
\[
  C_{i,k} \;=\; A_R \cdot \mathrm{NTT}(c_{i,k}) + B_R \cdot \mathrm{NTT}(r_{i,k})
\]
for $k = 0, \dots, t-1$, where $A_R, B_R$ are public matrices derived from \texttt{"pulsar.reshare.A.v1"} and \texttt{"pulsar.reshare.B.v1"} via BLAKE3-XOF $\to$ KeyedPRNG $\to$ UniformSampler. (These tags are distinct from the \texttt{dkg2} commit tags, so a DKG commit cannot be replayed as a reshare commit.) Each recipient $j$ verifies
\[
  A_R \cdot \mathrm{NTT}(\mathit{share}_{i \to j}) + B_R \cdot \mathrm{NTT}(\mathit{blind}_{i \to j}) \;=\; \sum_{k=0}^{t-1} \beta_j^k \cdot C_{i,k}
\]
in $R_q^M$ as exact equality. Hiding follows from MLWE on $B_R$; binding (over the small-norm cone) from MSIS on $[A_R \mid B_R]$. The argument is structurally identical to the \texttt{dkg2} commit scheme analyzed in \S\ref{sec:pedersen-dkg}.

\paragraph{Complaints and disqualification} (\texttt{complaint.go}). On a failed verification --- \textsc{BadDelivery}, \textsc{Equivocation}, \textsc{Missing}, or \textsc{MalformedCommit} --- the recipient broadcasts an Ed25519-signed \texttt{Complaint} naming the misbehaving sender and the failing slot. Sender $i$ is disqualified iff $\geq t_{\mathrm{old}} - 1$ distinct complainers signed valid complaints against $i$. The threshold $t_{\mathrm{old}} - 1$ exceeds the static-corruption bound $t_{\mathrm{old}} - 1$ by exactly one in the worst-case dishonest coalition, so no honest sender can be disqualified by a maximally adversarial set. After disqualification settles, every honest party computes the same surviving quorum $\mathcal{Q}' = \mathcal{Q} \setminus \mathit{disqualified}$ and the same $\lambda^{\mathcal{Q}'}$ coefficients. If $|\mathcal{Q}'| < t_{\mathrm{old}}$, the round fails and the chain stays at the old epoch.

\paragraph{Slashing evidence.} Each disqualifying \texttt{Complaint}, plus the equivocation pair $(C_a, C_b)$ where applicable, is admissible as slashing evidence at the consensus layer. The slasher re-runs \texttt{VerifyShareAgainstCommits} or \texttt{CommitDigest} equality under the recorded transcript hash; successful re-verification triggers the slashing transaction. The on-chain wire format is documented in \texttt{quasar\_integration.go}.

\subsection{Activation circuit-breaker}
\label{sec:resharing-activation}

Resharing is meaningless without a hard go/no-go switch at the chain level. If the new committee cannot collectively sign under the \emph{unchanged} group public key, the shares it holds do not actually reconstruct the master secret, and rotating to those shares would brick the chain. The \emph{activation certificate} is the proof that the new committee \emph{can} sign --- and therefore the chain accepts the new epoch.

The activation message (canonical bytes signed by the new committee under the unchanged group public key):
\begin{verbatim}
"QUASAR-PULSAR-ACTIVATE-v1"
    || transcript_hash         (32 bytes, TranscriptInputs.Hash())
    || reshare_transcript_hash (32 bytes, exchange digest:
                                 commits, complaints, disqualifications)
\end{verbatim}
The \texttt{transcript\_hash} field is a structural binding to the resharing parameters. The \texttt{reshare\_transcript\_hash} field is the BLAKE3 digest of the actual exchange under personalization \texttt{"pulsar.reshare.exchange-transcript.v1"}: every committed value, every complaint, every disqualification, in canonical order. Two honest validators with the same view of the exchange compute the same 32-byte hash.

\paragraph{Acceptance rule (consensus layer).}
\begin{enumerate}
\item The chain extracts the new $\textsc{EpochShareState}$ from the resharing batch.
\item The new committee runs the Sign protocol on the activation message.
\item The threshold-aggregated signature MUST verify under the unchanged group public key.
\item If verification succeeds, the chain commits the new $\textsc{EpochShareState}$, forwards activation evidence to the slashing module, and erases the old shares (\S\ref{sec:resharing-erasure}).
\item If verification fails, the chain stays at the old epoch. The old committee continues signing. A new resharing attempt is scheduled at the next epoch boundary.
\end{enumerate}

\subsection{KeyShare regeneration}
\label{sec:resharing-keyshare}

The Pulsar and Ringtail signing protocols consume a per-party \texttt{KeyShare} struct that carries more than just the bare Shamir share:
\[
  \texttt{KeyShare} = (\mathit{Index},\, \mathit{SkShare},\, \mathit{Seeds},\, \mathit{MACKeys},\, \Lambda,\, \mathit{GroupKey}).
\]
\textsc{Reshare} and \textsc{Refresh} only update $\mathit{SkShare}$. The remaining four fields must be regenerated for a complete \texttt{KeyShare} that can drive the signing protocol:

\begin{itemize}
\item $\Lambda$ (Lagrange coefficient): deterministic from the new committee's evaluation point set, recomputed via \texttt{primitives.ComputeLagrangeCoefficients} and brought into NTT-Montgomery form.

\item $\mathit{Seeds}_{ij}$ (per-pair PRF seed): derived from $\mathit{auth\_kex}_{ij}$ under the BLAKE3 personalization \texttt{"pulsar.reshare.prf-seed.v1"} with explicit binding to $(\mathit{chain\_id}, \mathit{group\_id}, \mathit{new\_epoch\_id}, i, j)$. The pair ordering $(i, j)$ is canonicalized (smaller-id first) so both endpoints derive the same value.

\item $\mathit{MACKeys}_{ij}$ (per-pair MAC key): same construction as $\mathit{Seeds}$, but under personalization \texttt{"pulsar.reshare.mac-key.v1"}. The distinct tag enforces $\mathit{Seeds}_{ij} \neq \mathit{MACKeys}_{ij}$ even though both are KDF outputs of the same $\mathit{auth\_kex}_{ij}$.

\item $\mathit{GroupKey}$: pointer copy of the persistent $(A, \tilde b)$ --- \emph{unchanged} across resharing.
\end{itemize}

The regeneration is implemented in \texttt{keyshare.go} as \texttt{PartyKey\-Share\-From\-Share}. The integration test \texttt{TestFullReshareThenKeyShareThenSign} drives the full chain Reshare $\to$ \texttt{PartyKey\-Share\-From\-Share} $\to$ Sign and confirms that the resulting signature verifies against the original (unchanged) $\tilde b$.

\subsection{Forced erasure}
\label{sec:resharing-erasure}

After the activation cert is accepted, every old share MUST be erased. The package provides \texttt{EraseShare} as a reasonable-effort overwrite of every \texttt{uint64} coefficient in every level of every polynomial in the share. Production deployments SHOULD additionally:
\begin{enumerate}
\item Pin the share allocation with \texttt{mlock} (Linux) or equivalent so paged-out swap cannot leak it after the fact.
\item Call \texttt{EraseShare} inline at the activation moment; do not retain stale references to the old share.
\item Optionally trigger \texttt{runtime.GC()} to free non-pinned references; pinned memory is unaffected.
\end{enumerate}
Failure to erase undermines the proactive-security guarantee: a long-running adversary that compromises a node \emph{after} resharing-activation gains the old share if it remains in memory, and an old-share leak combined with $< t_{\mathrm{new}}$ new shares does \emph{not} reconstruct the master secret --- but it does pin one position of the original Shamir polynomial, eroding the threshold security margin in a way that aggregates over many epochs.

\subsection{Quasar consensus integration}
\label{sec:resharing-quasar}

The Quasar epoch manager (\texttt{consensus/protocol/quasar/epoch.go}) currently regenerates a \emph{fresh} group public key at every validator-set rotation by calling \texttt{ringtailThreshold.GenerateKeys}. Under \textsc{Reshare}, the group public key is invariant for the entire group lineage; the rotation method is renamed for clarity:

\begin{center}
\begin{tabular}{ll}
\textbf{Old name} & \textbf{New name} \\ \hline
\texttt{RotateEpochKeys} & \texttt{ReshareEpoch} \\
\texttt{EpochKeys} & \texttt{EpochShareState} \\
\texttt{RotateEpoch(newVals)} calling \texttt{GenerateKeys} & \texttt{ReshareEpoch(newVals)} calling \texttt{Reshare(\ldots)} \\
GroupKey changes per epoch & GroupKey persistent for group lineage \\
\end{tabular}
\end{center}

\noindent
The forward-compat sketch in \texttt{consensus/protocol/quasar/reshare\_epoch.go} declares \texttt{ReshareEpoch} as an alias for \texttt{RotateEpoch} pending the body change, and exposes \texttt{ActivationMessage\-Personalization} plus the canonical transcript-input shape. The full integration checklist (renaming, transport reuse, complaint workflow, slashing evidence ingestion) is documented in \texttt{pulsar/reshare/quasar\_integration.go}.

\subsection{KAT cross-implementation contract}
\label{sec:resharing-kat}

The Go canonical at \texttt{pulsar/reshare/reshare.go} is byte-equivalent to the C++ port at \texttt{luxcpp/crypto/pulsar/reshare/} under a deterministic SHA-256 counter-mode PRNG. The KAT oracle at \texttt{pulsar/cmd/reshare\_oracle/main.go} emits 22 entries to \texttt{luxcpp/crypto/pulsar/test/kat/reshare\_kat.json}: 16 \textsc{Reshare} entries and 6 \textsc{Refresh} entries.

The 16 \textsc{Reshare} entries cover:
\begin{itemize}
\item Threshold tightening: $(2, 3) \to (3, 5)$, $(3, 5) \to (5, 7)$, $(5, 7) \to (7, 11)$
\item Threshold loosening: $(5, 7) \to (3, 5)$, $(5, 7) \to (2, 3)$
\item Same-threshold rotation: $(3, 5) \to (3, 5)$ disjoint, $(5, 7) \to (5, 7)$ disjoint
\item Committee growth: $(2, 3) \to (5, 9)$, $(2, 3) \to (2, 9)$
\item Partial overlap: $(3, 5) \to (3, 5)$ with shared and fresh IDs
\item Edge: $n_{\mathrm{new}} = t_{\mathrm{new}} + 1$ (smallest valid new committee)
\end{itemize}

The 6 \textsc{Refresh} entries cover $t \in \{1, 2, 3, 5, 7\}$ across canonical and relabelled committee IDs, including the degenerate $t = 1$ case (where \textsc{Refresh} must be the identity). Each entry pins (i) the planted secret $s$, (ii) the deterministic seed for the old Shamir polynomial coefficients, (iii) the seed for the resharing/refresh PRNG, (iv) the byte-exact old shares, (v) the byte-exact new shares, and (vi) per-share SHA-256 commitments. The C++ test reads the JSON, replays both primitives under the pinned PRNG seeds, and asserts byte-equality on all 22 entries.

\subsection{Performance}

The dominant cost of \textsc{Reshare} is dealing: each old party in $\mathcal{Q}$ samples $(t_{\mathrm{new}} - 1)$ polynomials in $R_q^{n_{\textrm{vec}}}$ ($\approx t_{\mathrm{new}} \cdot n_{\textrm{vec}} \cdot \varphi$ uniform $\mathbb{F}_q$ samples, totaling $\approx 7 \cdot 256 \cdot t_{\mathrm{new}}$ samples per party for the production parameters $n_{\textrm{vec}} = 7$, $\varphi = 256$) and evaluates the resulting polynomial at $|\mathcal{N}|$ points. Each evaluation is a Horner sweep over $t_{\mathrm{new}}$ coefficients, dominated by $t_{\mathrm{new}} \cdot n_{\textrm{vec}} \cdot \varphi$ modular multiplications.

For the canonical Quasar-21 deployment ($t_{\mathrm{new}} = 14$, $|\mathcal{N}| = 21$), each party performs $\approx 14 \cdot 21 \cdot 7 \cdot 256 \approx 5 \cdot 10^5$ modular multiplications. On a single Apple Silicon M2 Ultra core this completes in well under $10$~ms, dominated by the BLAKE3-XOF stream for the random coefficients. The wire cost is $|\mathcal{Q}| \cdot |\mathcal{N}| \cdot n_{\textrm{vec}} \cdot \varphi \cdot \lceil \log_2 q \rceil / 8 \approx 14 \cdot 21 \cdot 7 \cdot 256 \cdot 6 = 3.16$~MiB total across the committee, or $\approx 150$~KiB per ordered pair. The latency is dominated by the slowest pairwise channel, not by computation.

In comparison, a fresh trusted-dealer Gen at the same parameters costs the same $\approx 5 \cdot 10^5$ modular multiplications per share plus an additional $A \cdot s$ multiplication ($m_{\textrm{pk}} \cdot n_{\textrm{vec}} \cdot \varphi^2 / 2 \approx 1.8 \cdot 10^6$ multiplications under NTT) and the discrete-Gaussian sampling for $s, e$. The crucial cost difference is not asymptotic but architectural: \textsc{Reshare} preserves the public key, avoids a single dealer holding $s$, and runs entirely on operational infrastructure --- no offline ceremony required.
